% Catalogo de figuras propuestas para el capitulo 2 de la tesis.
% Compilar: pdflatex propuesta-figuras.tex (dos pasadas). No va en la tesis,
% es para evaluar antes de mover cada figura a escritos/04_antecedentes.
\documentclass[a4paper,10pt]{article}
\usepackage[utf8]{inputenc}
\usepackage[spanish,es-tabla]{babel}
\usepackage[margin=2.2cm]{geometry}
\usepackage{amsmath,amssymb}
\usepackage{xcolor}
\usepackage{tikz}
\usetikzlibrary{arrows.meta,positioning,fit,backgrounds,calc,shapes.geometric}
\usepackage{pgfplots}
\pgfplotsset{compat=1.18}
\usepackage{subcaption}
\usepackage{caption}
\captionsetup{font=small,labelfont=bf}
\usepackage{float}
\floatplacement{figure}{H} % en el catalogo cada figura va donde se escribe

% Paleta tomada del manual del sistema (docs/manual). En la tesis se usan
% sobre todo azul y naranja, con el gris de siempre; el resto solo donde hace
% falta distinguir mas de dos cosas.
\definecolor{azul}{HTML}{4C78A8}
\definecolor{naranja}{HTML}{F58518}
\definecolor{verde}{HTML}{54A24B}
\definecolor{rojo}{HTML}{E45756}
\definecolor{violeta}{HTML}{B279A2}
\definecolor{gris}{HTML}{8C8C8C}

\newcommand{\donde}[1]{\par\noindent\textbf{Dónde va.} #1\par\smallskip}
\newcommand{\pie}[1]{\par\noindent\textbf{Pie propuesto.} #1\par\smallskip}
\newcommand{\nota}[1]{\par\noindent\textbf{Nota.} #1\par\smallskip}

\title{Figuras propuestas para el Capítulo 2}
\author{Catálogo para evaluar, no forma parte de la tesis}
\date{Octubre de 2026}

\begin{document}
\maketitle

\noindent
Cada entrada dice si la figura es nueva o reemplaza una existente, en qué
sección y después de qué párrafo iría, y el pie propuesto. Los colores son
los del manual del sistema, azul \tikz\fill[azul] (0,0) rectangle (0.8em,0.8em);
y naranja \tikz\fill[naranja] (0,0) rectangle (0.8em,0.8em); como acentos,
gris para lo que no se quiere destacar, y verde \tikz\fill[verde] (0,0) rectangle (0.8em,0.8em);
o violeta \tikz\fill[violeta] (0,0) rectangle (0.8em,0.8em); solo cuando hay
que distinguir más de dos cosas. La numeración de secciones es la nueva, con
2.3.4 otras decisiones, 2.4 perfil, 2.5 asignación y 2.6 afinidad.

\tableofcontents

% =====================================================================
\section{Nuevas}
% =====================================================================

\subsection{N1. Programa lineal y relajación entera}

\donde{Nueva, en 2.7 Formulación matemática, después del párrafo que presenta
la programación lineal y el símplex, antes de la Ecuación del problema de
asignación. El párrafo la referenciaría dos veces, en ``el mínimo existe y
hay al menos un vértice que lo alcanza'' y en ``los vértices del poliedro
relajado ya no tienen por qué ser enteros''.}

\begin{figure}[H]
    \centering
    \begin{subfigure}{0.48\textwidth}
        \centering
        \begin{tikzpicture}[x=1.05cm, y=1.05cm]
            % region factible: x>=0, y>=0, x+2y<=7, 3x+y<=9
            \fill[azul!12] (0,0) -- (3,0) -- (2.2,2.4) -- (0,3.5) -- cycle;
            \draw[gris, thin] (-0.3,3.65) -- (3.6,1.7) node[pos=0.78, below, font=\scriptsize, black!60, sloped] {$x + 2y \le 7$};
            \draw[gris, thin] (3.3,-0.9) -- (1.7,3.9) node[pos=0.9, right, font=\scriptsize, black!60] {$3x + y \le 9$};
            \draw[azul, thick] (0,0) -- (3,0) -- (2.2,2.4) -- (0,3.5) -- cycle;
            % ejes
            \draw[-{Stealth}, black!70] (-0.3,0) -- (3.9,0) node[below, font=\scriptsize] {$x$};
            \draw[-{Stealth}, black!70] (0,-0.3) -- (0,4.1) node[left, font=\scriptsize] {$y$};
            % curvas de nivel del objetivo x+3y = c, para c = 4,5; 7,5 y 10,5
            \begin{scope}
                \clip (-0.2,-0.2) rectangle (3.8,4.0);
                \draw[naranja, dashed, thin] (0,1.5) -- (3.8,0.2333);
                \draw[naranja, dashed, thin] (0,2.5) -- (3.8,1.2333);
                \draw[naranja, dashed, thin] (0,3.5) -- (3.8,2.2333);
            \end{scope}
            \node[font=\scriptsize, naranja, anchor=west] at (2.55,2.95) {$x + 3y = 10{,}5$};
            \node[font=\scriptsize, naranja, anchor=west] at (2.95,1.55) {$x + 3y = 7{,}5$};
            % direccion de mejora
            \draw[-{Stealth}, naranja, thick] (1.0,0.4) -- ++(0.32,0.96) node[right, font=\scriptsize, pos=0.7] {mejora};
            % vertices
            \foreach \p in {(0,0),(3,0),(2.2,2.4)} \fill[azul] \p circle (0.055);
            \fill[naranja] (0,3.5) circle (0.08);
            \node[font=\scriptsize, anchor=west] at (0.14,3.72) {óptimo $(0;\,3{,}5)$};
        \end{tikzpicture}
        \caption{El programa lineal}
    \end{subfigure}
    \hfill
    \begin{subfigure}{0.48\textwidth}
        \centering
        \begin{tikzpicture}[x=1.05cm, y=1.05cm]
            \fill[azul!12] (0,0) -- (3,0) -- (2.2,2.4) -- (0,3.5) -- cycle;
            \draw[azul, thick] (0,0) -- (3,0) -- (2.2,2.4) -- (0,3.5) -- cycle;
            \draw[-{Stealth}, black!70] (-0.3,0) -- (3.9,0) node[below, font=\scriptsize] {$x$};
            \draw[-{Stealth}, black!70] (0,-0.3) -- (0,4.1) node[left, font=\scriptsize] {$y$};
            % puntos enteros factibles
            \foreach \x in {0,...,3} \foreach \y in {0,...,3} {
                \pgfmathparse{(\x+2*\y<=7) && (3*\x+\y<=9) ? 1 : 0}
                \ifnum\pgfmathresult=1 \fill[black!55] (\x,\y) circle (0.05); \fi
            }
            % optimo relajado (hueco) y mejor entero
            \draw[naranja, thick, dashed, thin] (-0.3,3.6) -- (3.6,2.3);
            \draw[naranja, thick] (0,3.5) circle (0.09);
            \node[font=\scriptsize, anchor=west, naranja] at (0.14,3.75) {relajado, $10{,}5$};
            \fill[naranja] (1,3) circle (0.09);
            \node[font=\scriptsize, anchor=west] at (1.14,3.12) {mejor entero, $10$};
            \draw[black!60, thick] (0,3) circle (0.09);
            \node[font=\scriptsize, anchor=east, black!60] at (-0.14,2.95) {redondeo, $9$};
        \end{tikzpicture}
        \caption{Con variables enteras}
    \end{subfigure}
\end{figure}

\pie{Un programa lineal en dos variables y su versión entera. En (a) la región
factible es el polígono que recortan las restricciones, las líneas punteadas
son curvas de nivel del objetivo $x + 3y$, la flecha indica hacia dónde
mejora, y el máximo se alcanza en un vértice. En (b) solo valen los puntos de
coordenadas enteras, el vértice óptimo de la relajación no es entero, el
mejor punto entero no es un vértice del polígono y tampoco es el redondeo del
óptimo relajado.}

\nota{La instancia está elegida para que el redondeo falle, $(0;\,3)$ vale 9
y $(1;\,3)$ vale 10. En la tesis el problema se escribe como minimización,
pero para el dibujo un máximo se lee igual.}

\clearpage
\subsection{N2. Planta del depósito con el flujo de las seis etapas}

\donde{Reemplaza la Figura 2.2 actual (las seis cajitas en fila), en 2.3.1
Dónde se decide el costo, después del párrafo que define las seis etapas. El
texto ya la referencia (``que la Figura 2.2 ordena''), así que no hay que
tocar prosa.}

\begin{figure}[H]
    \centering
    \begin{tikzpicture}[x=1cm, y=1cm, font=\scriptsize]
        % edificio
        \draw[black!70, thick] (0,0) rectangle (11,6.2);
        % dock: puertas abajo
        \foreach \x in {1.2,2.2,3.2} \draw[fill=naranja!25, draw=naranja!70!black] (\x-0.35,-0.02) rectangle ++(0.7,0.18);
        \foreach \x in {7.8,8.8,9.8} \draw[fill=azul!25, draw=azul!70!black] (\x-0.35,-0.02) rectangle ++(0.7,0.18);
        \node[anchor=north] at (2.2,-0.1) {\emph{dock} de recepción};
        \node[anchor=north] at (8.8,-0.1) {\emph{dock} de despacho};
        % zona de recepcion y zona de empaque
        \draw[fill=naranja!10, draw=naranja!60!black, dashed] (0.4,0.4) rectangle (4.0,1.5);
        \node at (2.2,0.95) {recepción};
        \draw[fill=azul!10, draw=azul!60!black, dashed] (7.0,0.4) rectangle (10.6,1.5);
        \node at (8.8,0.95) {empaque y despacho};
        % estanterias: 6 pasillos verticales
        \foreach \i in {0,...,6} \draw[fill=gray!25, draw=black!60] (1.0+\i*1.5, 2.3) rectangle ++(0.5, 3.4);
        \node at (5.5,5.95) {almacenamiento};
        % flechas del flujo
        \draw[-{Stealth}, very thick, naranja] (2.2,0.18) -- (2.2,0.4);
        \draw[-{Stealth}, very thick, naranja] (2.2,1.5) -- (2.2,2.0) -- (3.25,2.0) -- (3.25,3.6);
        \node[naranja!70!black, anchor=west] at (2.35,1.8) {\emph{put-away}};
        % picking: recorrido en S por dos pasillos
        \draw[-{Stealth}, very thick, azul] (5.5,1.5) -- (5.5,2.0) -- (6.25,2.0) -- (6.25,5.9) -- (7.75,5.9) -- (7.75,2.0) -- (8.8,2.0) -- (8.8,1.5);
        \fill[naranja] (6.1,3.1) circle (0.08);
        \fill[naranja] (6.4,4.6) circle (0.08);
        \fill[naranja] (7.9,3.9) circle (0.08);
        \node[azul!70!black, anchor=east] at (6.1,2.45) {\emph{picking}};
        \draw[-{Stealth}, very thick, azul] (8.8,0.4) -- (8.8,0.18);
        % numeros de etapa
        \foreach \p/\n in {(2.2,0.3)/1, (2.6,2.2)/2, (4.0,4.0)/3, (7.0,5.4)/4, (8.3,1.2)/5, (9.4,0.3)/6}
            \node[circle, fill=white, draw=black!70, inner sep=1pt, minimum size=3.2mm] at \p {\tiny\n};
    \end{tikzpicture}
\end{figure}

\pie{Planta esquemática de un depósito con el camino de la mercadería.
Entra por el \emph{dock} de recepción (1), el \emph{put-away} la lleva a una
ubicación (2) donde queda almacenada (3), el \emph{picking} la recolecta
recorriendo los pasillos (4) y la deja en empaque (5), y sale por el
\emph{dock} de despacho (6). Los rectángulos grises son las estanterías, los
puntos naranjas las ubicaciones que visita un recorrido de \emph{picking}.}

\nota{Es un flujo pasante, recepción de un lado y despacho del otro, para que
las seis etapas se lean de izquierda a derecha. En 2.3.2 se explica que
también puede ser en U.}

\subsection{N3. Quién se mueve}

\donde{Nueva, en 2.2 La digitalización y el \emph{warehouse} moderno, al lado
del párrafo ``La industria ensayó dos respuestas a esa presión'', que presenta
\emph{parts-to-picker} y \emph{picker-to-parts} y fija el alcance del
trabajo.}

\begin{figure}[H]
    \centering
    \begin{subfigure}{0.48\textwidth}
        \centering
        \begin{tikzpicture}[x=1cm, y=1cm, font=\scriptsize]
            \foreach \i in {0,...,4} \draw[fill=gray!25, draw=black!60] (\i*1.0, 0.6) rectangle ++(0.4, 2.6);
            % operario que camina
            \draw[-{Stealth}, very thick, azul] (0.2,0.2) -- (1.7,0.2) -- (1.7,2.2) -- (3.7,2.2) -- (3.7,0.2) -- (4.4,0.2);
            \fill[naranja] (1.45,2.2) circle (0.08);
            \fill[naranja] (3.95,1.3) circle (0.08);
            \node[circle, fill=azul, inner sep=0pt, minimum size=2.4mm] at (0.2,0.2) {};
            \node[anchor=north] at (0.2,0.05) {operario};
            \node[anchor=north] at (2.2,-0.35) {el operario camina hasta el producto};
        \end{tikzpicture}
        \caption{\emph{Picker-to-parts}}
    \end{subfigure}
    \hfill
    \begin{subfigure}{0.48\textwidth}
        \centering
        \begin{tikzpicture}[x=1cm, y=1cm, font=\scriptsize]
            % estanterias moviles
            \foreach \i in {0,...,3} \draw[fill=gray!25, draw=black!60] (0.3+\i*1.0, 1.6) rectangle ++(0.5, 1.6);
            % estacion fija
            \draw[fill=azul!15, draw=azul!70!black] (3.6,0.0) rectangle (4.6,0.9);
            \node[circle, fill=azul, inner sep=0pt, minimum size=2.4mm] at (4.1,0.45) {};
            \node[anchor=north] at (4.1,-0.05) {estación};
            % una estanteria viaja hasta la estacion
            \draw[fill=gray!25, draw=black!60] (2.6,0.1) rectangle ++(0.5,0.9);
            \fill[naranja] (2.85,0.55) circle (0.08);
            \draw[-{Stealth}, very thick, naranja] (1.55,1.55) -- (2.4,0.75);
            \node[anchor=north] at (2.2,-0.35) {el producto viaja hasta el operario};
        \end{tikzpicture}
        \caption{\emph{Parts-to-picker}}
    \end{subfigure}
\end{figure}

\pie{Las dos formas de juntar un pedido. En (a) el operario recorre los
pasillos y se detiene en las ubicaciones que necesita (puntos naranjas). En
(b) el operario trabaja quieto en una estación y las estanterías, o
contenedores, se le acercan. Este trabajo estudia el primer esquema.}

\clearpage
\subsection{N4. Zonas ABC, en franjas o en pasillos enteros}

\donde{Nueva, en 2.5.4 Clases y restricciones, junto a la oración ``Las zonas
pueden dibujarse de dos maneras, con las A ocupando el tramo de cada pasillo
más cercano al corredor frontal, o con las A ocupando enteros los pasillos
más cercanos al \emph{dock}''.}

\begin{figure}[H]
    \centering
    \begin{subfigure}{0.48\textwidth}
        \centering
        \begin{tikzpicture}[x=1cm, y=1cm, font=\scriptsize]
            \foreach \i in {0,...,5} {
                \fill[azul!75] (\i*0.7, 0.4) rectangle ++(0.3, 0.9);
                \fill[azul!40] (\i*0.7, 1.3) rectangle ++(0.3, 0.9);
                \fill[azul!15] (\i*0.7, 2.2) rectangle ++(0.3, 1.0);
                \draw[black!60] (\i*0.7, 0.4) rectangle ++(0.3, 2.8);
            }
            \fill[naranja] (1.9, 0.05) circle (0.07);
            \node[anchor=north] at (1.9,-0.05) {\emph{dock}};
            \node[anchor=west] at (4.3,0.85) {A};
            \node[anchor=west] at (4.3,1.75) {B};
            \node[anchor=west] at (4.3,2.7) {C};
        \end{tikzpicture}
        \caption{Por franjas, a lo largo de cada pasillo}
    \end{subfigure}
    \hfill
    \begin{subfigure}{0.48\textwidth}
        \centering
        \begin{tikzpicture}[x=1cm, y=1cm, font=\scriptsize]
            \foreach \i/\c in {0/azul!15, 1/azul!40, 2/azul!75, 3/azul!75, 4/azul!40, 5/azul!15} {
                \fill[\c] (\i*0.7, 0.4) rectangle ++(0.3, 2.8);
                \draw[black!60] (\i*0.7, 0.4) rectangle ++(0.3, 2.8);
            }
            \fill[naranja] (1.9, 0.05) circle (0.07);
            \node[anchor=north] at (1.9,-0.05) {\emph{dock}};
            \node[anchor=south] at (1.9,3.3) {A};
            \node[anchor=south] at (0.85,3.3) {B};
            \node[anchor=south] at (3.0,3.3) {B};
            \node[anchor=south] at (0.15,3.3) {C};
            \node[anchor=south] at (3.65,3.3) {C};
        \end{tikzpicture}
        \caption{Por pasillos enteros}
    \end{subfigure}
\end{figure}

\pie{Dos maneras de dibujar las zonas de una clasificación ABC sobre el mismo
depósito, con el \emph{dock} abajo. El tono más oscuro es la clase A, la de
las SKU más pedidas. En (a) cada pasillo tiene su tramo A cerca del corredor
frontal, lo que conviene con políticas de ruteo que entran y salen por ese
corredor. En (b) los pasillos cercanos al \emph{dock} son enteros de clase A,
lo que conviene cuando cada pasillo que se entra se recorre completo.}

\subsection{N5. Agrupar primero, asignar después}

\donde{Nueva, en 2.8 Estado del arte, en el bloque ``Método de asignación'',
donde se dice que la norma es la descomposición en dos fases.}

\begin{figure}[H]
    \centering
    \begin{tikzpicture}[x=1cm, y=1cm, font=\scriptsize,
        sku/.style={circle, draw=black!60, fill=gray!20, inner sep=0pt, minimum size=4mm}]
        % fase 0: historial de pedidos
        \node[draw=black!60, fill=white, align=left, inner sep=4pt] (hist) at (0,1.5) {pedido 1: $A, B$\\pedido 2: $A, B, E$\\pedido 3: $C, D$\\pedido 4: $C, D$\\pedido 5: $E, F$\\$\cdots$};
        \node[anchor=south] at (0,2.75) {historial};
        \draw[-{Stealth}, thick, black!60] (1.4,1.5) -- (2.2,1.5) node[midway, above] {medir};
        % fase 1: grafo de afinidad y grupos
        \begin{scope}[shift={(3.4,0.2)}]
            \node[sku] (a) at (0,2.2) {$A$};
            \node[sku] (b) at (0.9,2.6) {$B$};
            \node[sku] (e) at (0.7,1.6) {$E$};
            \node[sku] (c) at (0.2,0.5) {$C$};
            \node[sku] (d) at (1.1,0.3) {$D$};
            \node[sku] (f) at (1.9,1.4) {$F$};
            \draw[azul, line width=1.8pt] (a) -- (b);
            \draw[azul, line width=1.0pt] (a) -- (e);
            \draw[azul, line width=0.8pt] (b) -- (e);
            \draw[azul, line width=1.8pt] (c) -- (d);
            \draw[azul, line width=0.6pt] (e) -- (f);
            \draw[naranja, dashed, thick] (0.45,2.1) ellipse (0.95 and 0.95);
            \draw[naranja, dashed, thick] (0.65,0.4) ellipse (0.85 and 0.5);
        \end{scope}
        \node[anchor=south] at (4.3,2.95) {grupos afines};
        \draw[-{Stealth}, thick, black!60] (5.7,1.5) -- (6.5,1.5) node[midway, above] {asignar};
        % fase 2: layout con grupos en zonas
        \begin{scope}[shift={(6.8,0.1)}]
            \foreach \i in {0,...,3} \draw[fill=gray!20, draw=black!60] (\i*0.8, 0.3) rectangle ++(0.35, 2.4);
            \fill[naranja!35] (0.0,0.3) rectangle (0.35,1.3);
            \node at (0.17,0.8) {\tiny $ABE$};
            \fill[naranja!35] (0.8,0.3) rectangle (1.15,1.0);
            \node at (0.97,0.65) {\tiny $CD$};
            \fill[naranja] (1.6,-0.05) circle (0.06);
            \node[anchor=north] at (1.6,-0.1) {\emph{dock}};
        \end{scope}
        \node[anchor=south] at (8.2,2.95) {asignación};
    \end{tikzpicture}
\end{figure}

\pie{El esquema de dos fases que domina la literatura. Del historial se mide
la afinidad entre SKU, representada por el grosor de cada línea, con ella se
arman grupos de SKU afines, y recién después se decide dónde va cada grupo
dentro del depósito, las zonas naranjas. Lo que se pierde al separar las dos
decisiones rara vez se mide.}

\clearpage
\subsection{N6. Perfil de actividad con datos reales}

\donde{Reemplaza la Figura 2.5 actual (las dos vistas esquemáticas), en 2.4
Perfil de actividad y demanda. El pie sigue siendo genérico y la lectura del
caso va en el texto, que ya dice ``la curva que aparece suele ser similar a la
de la Figura''.}

\IfFileExists{perfil-real/pareto.dat}{%
\begin{figure}[H]
    \centering
    \begin{subfigure}{0.48\textwidth}
        \centering
        \begin{tikzpicture}
            \begin{axis}[width=\linewidth, height=5cm, xlabel={Fracción de las SKU, ordenadas por picks}, ylabel={Fracción acumulada de picks}, xmin=0, xmax=1, ymin=0, ymax=1, xtick={0,0.2,0.5,1}, ytick={0,0.5,0.8,1}, tick label style={font=\scriptsize}, label style={font=\scriptsize}, grid=none]
                \addplot[very thick, azul] table[x=x, y=y] {perfil-real/pareto.dat};
                \draw[dashed, naranja] (axis cs:0.2,0) -- (axis cs:0.2,1) ;
            \end{axis}
        \end{tikzpicture}
        \caption{Curva de Pareto de los picks}
    \end{subfigure}
    \hfill
    \begin{subfigure}{0.48\textwidth}
        \centering
        \begin{tikzpicture}
            \begin{axis}[width=\linewidth, height=5cm, ybar, bar width=4pt, xlabel={Líneas por pedido}, ylabel={Fracción}, xmin=0.4, xmax=11.6, ymin=0, xtick={1,...,11}, xticklabels={1,2,3,4,5,6,7,8,9,10,11+}, tick label style={font=\scriptsize}, label style={font=\scriptsize}, legend style={font=\scriptsize, draw=none, at={(0.97,0.95)}, anchor=north east}]
                \addplot[fill=azul!45, draw=azul] table[x=k, y=frac_ordenes] {perfil-real/lineas.dat};
                \addplot[fill=naranja, draw=naranja!70!black] table[x=k, y=frac_picks] {perfil-real/lineas.dat};
                \legend{De los pedidos, De los picks}
            \end{axis}
        \end{tikzpicture}
        \caption{Distribución de líneas por pedido}
    \end{subfigure}
\end{figure}
\nota{Datos de \texttt{data/v2/medium}, 712.137 picks, 6.912 SKU con al menos un pick y 203.625 órdenes. En (a) el 20\,\% de las SKU más pedidas explica el 76\,\% de los picks, y 319 SKU, el 4,6\,\%, explican la mitad. En (b) hay un problema, las líneas por orden salen uniformes entre 2 y 5, sin órdenes de una línea ni de más de cinco, porque así las genera el simulador de ShipHero, no porque así sea la demanda. Ese panel no sirve para ilustrar lo que dice el texto, que los pedidos de una línea son la mayoría. Propongo usar (a) con datos del depósito y dejar (b) como esquema, o directamente dejar solo (a). Y el pie no debería decir ``real'' sin aclarar que las órdenes son simuladas.}

}{%
\nota{Los datos todavía se están calculando a partir de
\texttt{data/v2}. Cuando estén, esta figura muestra la curva de Pareto real
de los picks por SKU y la distribución real de líneas por pedido del
depósito, con los mismos ejes que el esquema actual.}
}

\pie{Dos vistas del perfil de actividad de un depósito real. En (a), las SKU
ordenadas de la más pedida a la menos pedida y la fracción acumulada de
picks que explican, con la línea punteada en el 20\,\% de las SKU. En (b),
para cada tamaño de pedido, qué fracción de los pedidos tiene ese tamaño y
qué fracción de los picks generan.}

\subsection{N7. Mapa del capítulo}

\donde{Nueva, en la apertura del capítulo, después del párrafo que lo
presenta. Responde al pedido de Juan de un diagrama con el recorrido. Es
opcional, la apertura ya cuenta ese recorrido en prosa.}

\begin{figure}[H]
    \centering
    \begin{tikzpicture}[x=1cm, y=1cm, font=\scriptsize,
        caja/.style={draw=black!60, rounded corners=2pt, align=center, minimum height=11mm, text width=23mm, inner sep=3pt},
        ctx/.style={caja, fill=gray!12},
        nucleo/.style={caja, fill=azul!15, draw=azul!70!black},
        lit/.style={caja, fill=naranja!15, draw=naranja!70!black},
        flecha/.style={-{Stealth}, thick, black!55}]
        \node[ctx] (s1) at (0,0) {2.1 Evolución\\[1pt]\tiny depósito, consolidación, SKU};
        \node[ctx, right=3mm of s1] (s2) {2.2 Digitalización\\[1pt]\tiny WMS, lista de recolección, historial};
        \node[ctx, right=3mm of s2] (s3) {2.3 Operaciones\\[1pt]\tiny etapas, dock, layout, picking, batching, ruteo};
        \node[ctx, right=3mm of s3] (s4) {2.4 Perfil\\[1pt]\tiny picks, ABC, líneas por pedido};
        \node[nucleo, below=7mm of s4] (s5) {2.5 Asignación\\[1pt]\tiny dedicado, reglas, $c_k$, COI};
        \node[nucleo, left=3mm of s5] (s6) {2.6 Afinidad\\[1pt]\tiny $n_{ij}$, seis medidas, matriz};
        \node[nucleo, left=3mm of s6] (s7) {2.7 Formulación\\[1pt]\tiny PL, asignación, QAP, heurísticas};
        \node[lit, left=3mm of s7] (s8) {2.8 y 2.9 Literatura\\[1pt]\tiny condiciones, resultados, huecos};
        \draw[flecha] (s1) -- (s2);
        \draw[flecha] (s2) -- (s3);
        \draw[flecha] (s3) -- (s4);
        \draw[flecha] (s4) -- (s5);
        \draw[flecha] (s5) -- (s6);
        \draw[flecha] (s6) -- (s7);
        \draw[flecha] (s7) -- (s8);
    \end{tikzpicture}
\end{figure}

\pie{Recorrido del capítulo. Las cuatro primeras secciones dan el contexto,
qué es un depósito, qué registra y cómo opera, y qué dicen sus datos. Las
tres siguientes son el núcleo, cómo se asignan ubicaciones, cómo se mide la
afinidad y cómo se escribe el problema. Las dos últimas revisan la
literatura y lo que deja abierto. Debajo de cada una, los términos que
define.}

\clearpage
\subsection{N8. Por qué el término cuadrático rompe el orden}

\donde{Nueva, en 2.7 Formulación matemática, después del párrafo que
introduce el término cuadrático y antes del QAP. Es la figura más chica del
lote y explica en un vistazo por qué la suma lineal se resuelve ordenando y la
cuadrática no.}

\begin{figure}[H]
    \centering
    \begin{tikzpicture}[x=1cm, y=1cm, font=\scriptsize,
        ub/.style={draw=black!60, fill=gray!15, minimum size=7mm, inner sep=0pt}]
        % fila (a): solo rotacion
        \node[anchor=east] at (-0.3,0.35) {(a) solo rotación};
        \foreach \i/\s/\f in {0/A/10, 1/B/8, 2/C/6, 3/D/4, 4/E/3, 5/F/2} {
            \node[ub] (a\i) at (\i*0.9,0.35) {$\s$};
            \node at (\i*0.9,-0.15) {$f=\f$};
        }
        \node[anchor=west] at (5.0,0.35) {\emph{dock} a la izquierda, $c_1 < c_2 < \cdots < c_6$};
        % fila (b): con afinidad A-E
        \node[anchor=east] at (-0.3,-1.6) {(b) con afinidad $A$--$E$};
        \foreach \i/\s/\f in {0/A/10, 1/E/3, 2/B/8, 3/C/6, 4/D/4, 5/F/2} {
            \node[ub] (b\i) at (\i*0.9,-1.6) {$\s$};
            \node at (\i*0.9,-2.1) {$f=\f$};
        }
        \draw[naranja, very thick] (b0.north) to[bend left=45] node[above] {$a_{AE}$ alto} (b1.north);
        \draw[azul, very thick] (a0.north) to[bend left=30] node[above] {$a_{AE}$ alto} (a4.north);
        \node[anchor=west, text width=5cm, align=left] at (5.0,-1.7) {$E$ sube tres lugares aunque se pida menos que $B$, $C$ y $D$, porque el par $A$--$E$ paga el viaje entre ambas en cada pedido que las junta};
    \end{tikzpicture}
\end{figure}

\pie{Seis SKU en seis ubicaciones de un mismo pasillo, ordenadas por
distancia al \emph{dock}. En (a) la suma lineal se minimiza poniendo la más
pedida más cerca, y el orden por $f_i$ es óptimo. En (b) aparece una
afinidad alta entre $A$ y $E$, y el término cuadrático, que paga la distancia
entre ambas por cada pedido que las junta, puede preferir acercarlas aunque
eso rompa el orden por $f_i$. Cuál de los dos pesa más depende de los
valores, y por eso el problema ya no se resuelve ordenando.}

% =====================================================================
\clearpage
\section{Existentes, recoloreadas}
% =====================================================================

\noindent
Las figuras actuales son todas en gris. Acá están con la misma paleta que las
nuevas, estanterías en gris, recorridos en azul, extracciones y \emph{docks} en
naranja. El dibujo no cambia, solo el color.

\subsection{R1. Figura 2.1, la unidad de manipulación}

\donde{2.1 Evolución, misma posición.}

\begin{figure}[H]
    \centering
    \begin{tikzpicture}[x=1cm, y=1cm, every node/.style={font=\small}]
        \foreach \i in {0,1,2} \foreach \j in {0,1,2}
            \draw[fill=azul!18, draw=azul!70!black] (\i*0.63, \j*0.5) rectangle ++(0.58, 0.44);
        \draw[fill=gray!55, draw=black!70] (-0.08, -0.2) rectangle (1.92, -0.04);
        \node at (0.92, -0.65) {\emph{pallet}};
        \draw[-{Stealth}, thick, black!60] (2.25, 0.7) -- (2.95, 0.7);
        \foreach \i in {0,1} \foreach \j in {0,1}
            \draw[fill=azul!18, draw=azul!70!black] (3.25+\i*0.46, 0.24+\j*0.46) rectangle ++(0.42, 0.42);
        \draw[draw=black!70, thick] (3.2, 0.19) rectangle (4.22, 1.21);
        \node at (3.71, -0.65) {caja};
        \draw[-{Stealth}, thick, black!60] (4.55, 0.7) -- (5.25, 0.7);
        \foreach \i in {0,1,2} \draw[fill=azul!18, draw=azul!70!black] (5.55+\i*0.36, 0.54) rectangle ++(0.32, 0.32);
        \draw[draw=black!70, thick] (5.5, 0.49) rectangle (6.66, 0.91);
        \node at (6.08, -0.65) {paquete interno};
        \draw[-{Stealth}, thick, black!60] (7.0, 0.7) -- (7.7, 0.7);
        \draw[fill=azul!18, draw=azul!70!black] (7.95, 0.54) rectangle ++(0.32, 0.32);
        \node at (8.11, -0.65) {unidad};
    \end{tikzpicture}
\end{figure}

\subsection{R2. Figura 2.3, variantes de geometría}

\donde{2.3.2 La geometría del depósito, misma posición. El sombreado de
conveniencia pasa de negro translúcido a azul, los \emph{docks} a naranja.}

\begin{figure}[H]
    \centering
    \begin{subfigure}{0.48\textwidth}
        \centering
        \begin{tikzpicture}[x=1cm, y=1cm]
            \foreach \i in {0,...,9} \draw[fill=gray!15, draw=black!60] (\i*0.42, 0) rectangle ++(0.18, 2.6);
            \fill[azul, opacity=0.18] (1.10, 0) rectangle (2.86, 2.6);
            \fill[azul, opacity=0.22] (1.52, 0) rectangle (2.44, 2.6);
            \fill[naranja] (1.98, -0.22) circle (0.07);
            \fill[naranja] (1.98, 2.82) circle (0.07);
            \node[font=\scriptsize] at (1.98, -0.5) {recepción};
            \node[font=\scriptsize] at (1.98, 3.08) {despacho};
        \end{tikzpicture}
        \caption{Flujo pasante}
    \end{subfigure}
    \hfill
    \begin{subfigure}{0.48\textwidth}
        \centering
        \begin{tikzpicture}[x=1cm, y=1cm]
            \foreach \i in {0,...,9} \draw[fill=gray!15, draw=black!60] (\i*0.42, 0) rectangle ++(0.18, 2.6);
            \fill[azul, opacity=0.14] (1.98, 2.35) -- (0.15, 0) -- (3.81, 0) -- cycle;
            \fill[azul, opacity=0.18] (1.98, 1.5) -- (0.85, 0) -- (3.11, 0) -- cycle;
            \fill[azul, opacity=0.24] (1.98, 0.75) -- (1.45, 0) -- (2.51, 0) -- cycle;
            \fill[naranja] (1.98, -0.22) circle (0.07);
            \node[font=\scriptsize] at (1.98, -0.5) {recepción y despacho};
            \node[font=\scriptsize, opacity=0] at (1.98, 3.08) {despacho};
        \end{tikzpicture}
        \caption{Flujo en U}
    \end{subfigure}

    \vspace{0.8em}

    \begin{subfigure}{0.48\textwidth}
        \centering
        \begin{tikzpicture}[x=1cm, y=1cm]
            \foreach \i in {0,...,9} {
                \draw[fill=gray!15, draw=black!60] (\i*0.42, 1.5) rectangle ++(0.18, 1.1);
                \draw[fill=gray!15, draw=black!60] (\i*0.42, 0) rectangle ++(0.18, 1.1);
            }
            \draw[very thick, azul] (0.72, 0.3) -- (0.72, 1.3) -- (3.24, 1.3) -- (3.24, 2.25);
            \fill[naranja] (0.63, 0.16) rectangle ++(0.18, 0.18);
            \fill[naranja] (3.15, 2.2) rectangle ++(0.18, 0.18);
            \fill[naranja] (1.98, -0.22) circle (0.07);
            \node[font=\scriptsize] at (1.98, -0.5) {dock};
        \end{tikzpicture}
        \caption{Pasillo transversal}
    \end{subfigure}
    \hfill
    \begin{subfigure}{0.48\textwidth}
        \centering
        \begin{tikzpicture}[x=1cm, y=1cm]
            \begin{scope}
                \clip (0, 0) -- (1.68, 0) -- (0.08, 2.6) -- (0, 2.6) -- cycle;
                \foreach \j in {0,...,5} \draw[fill=gray!15, draw=black!60] (0, 0.06+\j*0.44) rectangle (1.9, 0.24+\j*0.44);
            \end{scope}
            \begin{scope}
                \clip (2.28, 0) -- (3.96, 0) -- (3.96, 2.6) -- (3.88, 2.6) -- cycle;
                \foreach \j in {0,...,5} \draw[fill=gray!15, draw=black!60] (2.0, 0.06+\j*0.44) rectangle (3.96, 0.24+\j*0.44);
            \end{scope}
            \begin{scope}
                \clip (1.98, 0.2) -- (0.5, 2.6) -- (3.46, 2.6) -- cycle;
                \foreach \k in {0,...,8} \draw[fill=gray!15, draw=black!60] (0.15+\k*0.44, 0.2) rectangle ++(0.18, 2.4);
            \end{scope}
            \fill[naranja] (1.98, -0.22) circle (0.07);
            \node[font=\scriptsize] at (1.98, -0.5) {dock};
        \end{tikzpicture}
        \caption{Espina de pescado}
    \end{subfigure}
\end{figure}

\clearpage
\subsection{R3. Figura 2.4, las dos políticas de ruteo}

\donde{2.3.3 Order picking, misma posición. Recorrido en azul, extracciones
en naranja.}

\begin{figure}[H]
    \centering
    \begin{subfigure}{0.48\textwidth}
        \centering
        \begin{tikzpicture}[x=1cm, y=1cm]
            \foreach \i in {0,...,6} \draw[fill=gray!25, draw=black!60] (\i*0.62, 0.35) rectangle ++(0.24, 2.3);
            \draw[very thick, azul, -{Stealth}]
                (0.43, 0.15) -- (0.43, 2.85) -- (1.05, 2.85) -- (1.05, 0.15) -- (2.91, 0.15) -- (2.91, 2.85) -- (3.53, 2.85) -- (3.53, 0.15);
            \draw[thick, azul!60, dashed] (3.53, 0.15) -- (3.53, 0.02) -- (0.43, 0.02) -- (0.43, 0.15);
            \foreach \p in {(0.36, 1.9), (1.12, 0.9), (2.84, 1.4), (3.60, 2.1)} \fill[naranja] \p circle (0.07);
            \fill[black!70] (0.43, -0.2) circle (0.06);
            \node[font=\scriptsize] at (0.43, -0.45) {inicio};
        \end{tikzpicture}
        \caption{Política en forma de $S$}
    \end{subfigure}
    \hfill
    \begin{subfigure}{0.48\textwidth}
        \centering
        \begin{tikzpicture}[x=1cm, y=1cm]
            \foreach \i in {0,...,6} \draw[fill=gray!25, draw=black!60] (\i*0.62, 0.35) rectangle ++(0.24, 2.3);
            \draw[very thick, azul, -{Stealth}]
                (0.43, 0.15) -- (0.43, 1.9) -- (0.43, 0.15) -- (1.05, 0.15) -- (1.05, 0.9) -- (1.05, 0.15) -- (2.91, 0.15) -- (2.91, 1.4) -- (2.91, 0.15) -- (3.53, 0.15) -- (3.53, 2.1) -- (3.53, 0.15);
            \draw[thick, azul!60, dashed] (3.53, 0.15) -- (3.53, 0.02) -- (0.43, 0.02) -- (0.43, 0.15);
            \foreach \p in {(0.36, 1.9), (1.12, 0.9), (2.84, 1.4), (3.60, 2.1)} \fill[naranja] \p circle (0.07);
            \fill[black!70] (0.43, -0.2) circle (0.06);
            \node[font=\scriptsize] at (0.43, -0.45) {inicio};
        \end{tikzpicture}
        \caption{Política de retorno}
    \end{subfigure}
\end{figure}

\subsection{R4. Figura 2.6, el diente de sierra}

\donde{2.5.1 Dedicado o compartido, misma posición. Stock en azul, promedio
en naranja.}

\begin{figure}[H]
    \centering
    \begin{tikzpicture}
        \begin{axis}[width=0.7\textwidth, height=4.2cm, xlabel={Tiempo}, ylabel={Stock en la ubicación}, xmin=0, xmax=3, ymin=0, ymax=1.15, xtick={0,1,2,3}, xticklabels={$0$,$T$,$2T$,$3T$}, ytick={0.5,1}, yticklabels={$Q/2$,$Q$}, tick label style={font=\scriptsize}, label style={font=\scriptsize}, axis lines=left]
            \addplot[very thick, azul] coordinates {(0,1) (1,0) (1,1) (2,0) (2,1) (3,0)};
            \addplot[dashed, naranja, thick, domain=0:3] {0.5};
        \end{axis}
    \end{tikzpicture}
\end{figure}

\subsection{R5. Figura 2.7, dos ubicaciones y su costo de acceso}

\donde{2.5.2 Las reglas de asignación, misma posición. El viaje corto en
azul, el largo en naranja.}

\begin{figure}[H]
    \centering
    \begin{tikzpicture}[x=1cm, y=1cm]
        \foreach \i in {0,...,7} \draw[fill=gray!25, draw=black!60] (\i*0.62, 0.5) rectangle ++(0.24, 2.4);
        \fill[black!70] (2.29, 0.1) circle (0.07);
        \node[font=\scriptsize] at (2.29, -0.2) {dock};
        \fill[azul] (2.22, 1.0) circle (0.08);
        \node[font=\scriptsize, anchor=west] at (2.3, 1.0) {$A$};
        \draw[thick, azul, -{Stealth}] (2.29, 0.1) -- (2.29, 1.0);
        \draw[thick, azul, -{Stealth}] (2.36, 1.0) -- (2.36, 0.1);
        \fill[naranja] (0.36, 2.6) circle (0.08);
        \node[font=\scriptsize, anchor=east] at (0.3, 2.6) {$B$};
        \draw[thick, naranja, dashed, -{Stealth}] (2.29, 0.1) -- (0.43, 0.3) -- (0.43, 2.6);
        \draw[thick, naranja, dashed, -{Stealth}] (0.5, 2.6) -- (0.5, 0.36) -- (2.29, 0.16);
    \end{tikzpicture}
\end{figure}

\subsection{R6. Figura 2.5, el perfil esquemático}

\donde{Solo si N6 no se usa. Es el esquema actual con barras azul y naranja en
vez de gris claro y gris oscuro.}

\begin{figure}[H]
    \centering
    \begin{subfigure}{0.48\textwidth}
        \centering
        \begin{tikzpicture}
            \begin{axis}[width=\linewidth, height=4.6cm, xlabel={Fracción de las SKU, ordenadas por picks}, ylabel={Fracción acumulada de picks}, xmin=0, xmax=1, ymin=0, ymax=1, xtick={0,0.2,0.5,1}, ytick={0,0.5,0.8,1}, tick label style={font=\scriptsize}, label style={font=\scriptsize}, grid=none]
                \addplot[very thick, azul, domain=0:1, samples=100] {(1-exp(-8*x))/(1-exp(-8))};
                \draw[dashed, naranja] (axis cs:0.2,0) -- (axis cs:0.2,0.8) -- (axis cs:0,0.8);
            \end{axis}
        \end{tikzpicture}
        \caption{Curva de Pareto de los picks}
    \end{subfigure}
    \hfill
    \begin{subfigure}{0.48\textwidth}
        \centering
        \begin{tikzpicture}
            \begin{axis}[width=\linewidth, height=4.6cm, ybar, bar width=4pt, xlabel={Líneas por pedido}, ylabel={Fracción}, xmin=0.4, xmax=6.6, ymin=0, ymax=0.7, xtick={1,...,6}, ytick={0,0.2,0.4,0.6}, tick label style={font=\scriptsize}, label style={font=\scriptsize}, legend style={font=\scriptsize, draw=none, at={(0.97,0.95)}, anchor=north east}]
                \addplot[fill=azul!45, draw=azul] coordinates {(1,0.60) (2,0.20) (3,0.10) (4,0.05) (5,0.03) (6,0.02)};
                \addplot[fill=naranja, draw=naranja!70!black] coordinates {(1,0.34) (2,0.23) (3,0.17) (4,0.11) (5,0.08) (6,0.07)};
                \legend{De los pedidos, De los picks}
            \end{axis}
        \end{tikzpicture}
        \caption{Distribución de líneas por pedido}
    \end{subfigure}
\end{figure}

% =====================================================================
\clearpage
\section{Lo que no está acá}
% =====================================================================

\begin{itemize}
    \item \textbf{La foto de un depósito real.} La pidió Juan y es la única
    que no puedo generar. Iría en 2.1 o en 2.2, con un pie que diga qué se ve,
    pasillos, estanterías, el \emph{dock}.
    \item \textbf{El grafo de afinidad.} Se sacó porque ni la literatura
    revisada ni el método propio usan herramientas de grafos. Si el método
    llega a usarlas, vuelve con datos reales.
    \item \textbf{La Tabla 2.1 de tipos de depósito.} Juan la marcó como
    fea. Es una tabla, no una figura, y se rehace aparte.
\end{itemize}

\end{document}
